%-------------------------
% CV – Elyas Chikhaoui (FR)
% Compatible Overleaf (pdflatex)
%-------------------------
\documentclass[a4paper,10pt]{article}

\usepackage[T1]{fontenc}
\usepackage[utf8]{inputenc}
\usepackage[french]{babel}
\usepackage[left=1.5cm, right=1.5cm, top=1.0cm, bottom=1.0cm]{geometry}
\usepackage{charter}            % police sobre ; supprimer cette ligne pour revenir à Computer Modern
\usepackage[scaled=0.92]{helvet}
\usepackage{hyperref}
\usepackage{enumitem}
\usepackage{titlesec}
\usepackage{array}
\usepackage{tabularx}
\usepackage{xcolor}
\usepackage{fontawesome5}
\usepackage{microtype}

\frenchsetup{StandardLists=true}

%── Informations de contact (à modifier ici uniquement) ───────────────────────
\newcommand{\myname}{Elyas Chikhaoui}
\newcommand{\myphone}{+33 7 81 86 77 30}
\newcommand{\mymail}{elyas.chikhaoui@telecom-paris.fr}
\newcommand{\mygithub}{how2222}
\newcommand{\mylinkedin}{elyaschikhaoui}
\newcommand{\mysite}{how2222.github.io/portfolio}

%── Couleurs ──────────────────────────────────────────────────────────────────
\definecolor{accent}{RGB}{0,0,0}
\definecolor{muted}{RGB}{110,110,110}

\hypersetup{hidelinks, pdftitle={CV - \myname}, pdfauthor={\myname}}

%── Sections ──────────────────────────────────────────────────────────────────
\titleformat{\section}{\large\bfseries\scshape\color{accent}}{}{0em}{}[{\color{accent}\titlerule[0.6pt]}]
\titlespacing{\section}{0pt}{7pt}{4pt}

%── Entrée de CV : #1 organisation, #2 poste, #3 dates, #4 contenu ────────────
\newcommand{\cventry}[4]{%
  \noindent
  \begin{tabularx}{\textwidth}{@{}X@{\hspace{6pt}}r@{}}
    \textbf{#1} & \textbf{#3}\\
    \textit{#2} & \\
  \end{tabularx}
  \par\vspace{-1pt}
  #4
  \vspace{5pt}
}

\setlist[itemize]{leftmargin=11pt,noitemsep,topsep=1pt,parsep=0pt,itemsep=1.5pt,label=\textbullet}
\setlength{\parindent}{0pt}
\setlength{\parskip}{0pt}
\pagestyle{empty}

%─────────────────────────────────────────────────────────────────────────────
\begin{document}

\begin{center}
  {\LARGE\bfseries\scshape \myname}\\[4pt]
  {\small
    \faMapMarker*\ Paris
    \quad\textcolor{muted}{|}\quad \faPhone\ \myphone
    \quad\textcolor{muted}{|}\quad \faEnvelope\ \href{mailto:\mymail}{\mymail}
  }\\[3pt]
  {\small
    \faGithub\ \href{https://github.com/\mygithub}{GitHub}
    \quad\textcolor{muted}{|}\quad \faLinkedin\ \href{https://www.linkedin.com/in/\mylinkedin}{linkedin.com/in/\mylinkedin}
    \quad\textcolor{muted}{|}\quad \faGlobe\ \href{https://\mysite}{\mysite}
  }\\[4pt]
  {\small\color{muted}
    Math\'ematiques appliqu\'ees \& Machine Learning \textbar{} T\'el\'ecom Paris -- Institut Polytechnique de Paris
  }\\[2pt]
  {\small\bfseries
    Stage de fin d'\'etudes (6 mois) -- recherche ML / data science -- disponible \`a partir de mars 2027
  }
\end{center}

\section{Exp\'erience professionnelle \& associative}

\cventry{INRIA Paris}%
  {Stagiaire de recherche -- Machine Learning \& r\'egression symbolique}{janv.~2026 -- juil.~2026}{%
  \begin{itemize}
    \item Conception et impl\'ementation de \textbf{DragonSR}, un framework de r\'egression symbolique formulant la d\'ecouverte d'\'equations comme l'optimisation de graphes acycliques dirig\'es (DAG), construit au-dessus de DRAGON (Keisler et al., \textit{JMLR} 2024) : repr\'esentation multi-canaux avec partage de sous-expressions, op\'erateurs de s\'election de features, agr\'egation des canaux par moindres carr\'es (parcimonieuse, nested, polynomiale-rationnelle), recherche par bandits (Mutant-UCB).
    \item Pipeline robuste au bruit : d\'ebruitage par processus gaussien, sous-\'echantillonnage stochastique, augmentation de variables et s\'election par XGBoost, objectif de corr\'elation et p\'enalit\'e de complexit\'e structurelle.
    \item Benchmark de 27 \'equations (Nguyen, 9 lois physiques, 9 indices de t\'el\'ed\'etection Sentinel-2 \`a 10+ variables) face \`a \textbf{PySR} (\'etat de l'art), sur 4 niveaux de bruit $\times$ 10 runs : DragonSR retrouve les formules jusqu'\`a 10~\% de bruit l\`a o\`u PySR n'en retrouve aucune (Hubble, gaz parfaits, Kepler, Leavitt, SAVI, EVI2\ldots).
    \item \textbf{Premier auteur} de l'article \textit{DragonSR: Directed Acyclic Graph Search for Equation Discovery} (avec J.~Keisler et A.~Charantonis), pr\'epublication arXiv \`a para\^itre ; code et r\'esultats publi\'es (GitHub, Zenodo).
  \end{itemize}
}

\cventry{Independent Commodity Intelligence Services (ICIS) -- LexisNexis Risk Solutions}%
  {Stagiaire analyste quantitatif -- march\'es du gaz}{juil.~2025 -- d\'ec.~2025}{%
  \begin{itemize}
    \item Extension d'un mod\`ele LP du r\'eseau gazier europ\'een (GAMS) : proposition et impl\'ementation d'une approche par fonctions triangulaires pour mod\'eliser le sendout GNL aux terminaux, rempla\c{c}ant une d\'echarge continue par un profil de d\'echargement des m\'ethaniers distribu\'e dans le temps.
    \item D\'eveloppement de 15+ graphiques analytiques sur le portail propri\'etaire ICIS (SQL) : spreads de localisation, courbes forward, utilisation des gazoducs ; pipeline de donn\'ees ing\'erant les r\'ealis\'es manquants via l'API ENTSOG pour alimenter des indicateurs de stress d'offre flexible ; carte interactive des flux et prix gaziers europ\'eens quotidiens et pr\'evisionnels, d\'evelopp\'ee de z\'ero (Plotly, Dash).
    \item Refonte et modularisation du pipeline de traitement des indisponibilit\'es sous un nouveau mod\`ele de donn\'ees : temps d'ex\'ecution par module r\'eduit de $\sim$3~h \`a $\sim$20~min (Python, SQL).
  \end{itemize}
}

\cventry{Sopra Steria}%
  {Stagiaire -- Digital Unit, division Aeroline}{juil.~2024 -- ao\^ut~2024}{%
  \begin{itemize}
    \item Contribution \`a \textit{Bhubble}, application web Angular de gestion d'\'equipes et de projets (int\'egration GitLab, sondages, veille technologique).
    \item Script Python d'automatisation de la migration de tickets (Tuleap $\to$ Spira) ; travail en m\'ethode Agile avec TypeScript, API REST et PostgreSQL.
  \end{itemize}
}

\cventry{T\'el\'ecom IA \& Unaite -- associations \'etudiantes d'IA}%
  {Pr\'esident \& co-fondateur}{juin~2024 -- juil.~2025}{%
  \begin{itemize}
    \item Pr\'esident de \textit{T\'el\'ecom IA} : organisation d'\'ev\'enements et d'ateliers techniques avec InstaDeep et Kyutai.
    \item Co-fondateur d'\textit{Unaite}, alliance nationale des associations \'etudiantes d'IA de 6 grandes \'ecoles, 100+ membres (T\'el\'ecom Paris, \'Ecole Polytechnique, ENPC, CentraleSup\'elec, HEC, ESSEC).
  \end{itemize}
}

\section{Formation}

\cventry{T\'el\'ecom Paris \textbar{} Institut Polytechnique de Paris}%
  {Dipl\^ome d'ing\'enieur -- Math\'ematiques appliqu\'ees}{sept.~2023 -- juin~2027}{%
  \begin{itemize}
    \item \textbf{Cours :} calcul stochastique, probabilit\'es en grande dimension, machine learning, informatique, sociologie.
    \item \textbf{Semestre d'\'echange \`a AGH (Cracovie, Pologne)}, automne 2026 : algorithmes stochastiques, traitement d'images et vision, programmation GPU (CUDA / OpenCL), mod\'elisation et optimisation avec GAMS, soft computing, semi-conducteurs pour le calcul neuromorphique.
  \end{itemize}
}

\cventry{Lyc\'ee Henri IV}%
  {Classes pr\'eparatoires MPSI \& MP}{sept.~2020 -- juil.~2023}{%
  \begin{itemize}
    \item Pr\'eparation intensive en math\'ematiques et physique aux concours d'entr\'ee des grandes \'ecoles.
  \end{itemize}
}

\section{Comp\'etences \& centres d'int\'er\^et}

\begin{tabularx}{\textwidth}{@{}>{\bfseries}p{2.1cm} X@{}}
  Langues     & Fran\c{c}ais (langue maternelle), Anglais (courant, Cambridge Linguaskill : C1), Polonais (conversationnel) \\[2pt]
  Techniques  & Python (PyTorch, scikit-learn, Pandas, NumPy, Plotly, Dash), SQL / PostgreSQL, GAMS, Git, \LaTeX{} \\[2pt]
  Int\'er\^ets   & Karat\'e (ceinture noire), Football \\
\end{tabularx}

\end{document}
